\documentclass[11pt]{article}

\usepackage[a4paper,margin=1in]{geometry}
\usepackage{amsmath,amssymb,amsthm,mathtools}
\usepackage{physics}
\usepackage{bbm}
\usepackage{hyperref}
\usepackage{enumitem}
\usepackage{mathrsfs}
\usepackage{bm}

\title{Unified Prime-Weighted Evolution Operator \texorpdfstring{$\Xi(t)$}{}\\
for Hydrogenic Spectra and Hypercomputational Number Fields}
\author{Tyler van Osdol, Martin Gibson \& Ryan O. van Gelder \\ \\ Citizen Gardens \\ \\ The Foundation of Multiplicity}
\date{April 3, 2026}

\begin{document}
\maketitle

\begin{abstract}
We present a unified prime-weighted, log-periodic evolution operator
$\Xi(t)$ acting simultaneously on (i) hydrogenic states expressed in
Hydrogen Universe Standard (H-US) units and (ii) number-theoretic
structures arising in the $\Phi$-Quantum Natural Number Singularity
($\Phi$-QNNS) framework. On the physical side, $\Xi(t)$ deforms
hydrogen $s$-state wavefunctions through a non-commutative product of
prime-indexed dilation operators, yielding log-radius ($\log\rho$)
spectral signatures with enhanced power at frequencies $\omega_p =
\log p$. On the arithmetic side, the same prime weights and
frequencies drive a gauge-theoretic recursion whose fixed points
realise the hierarchy $\mathbb{N} \subset \mathbb{Q} \subset
\mathbb{R} \subset \mathbb{C}$ as a stratified sheaf in a
$\Phi$-topos. We give a compact mathematical definition of the unified
$\Xi(t)$, outline its numerical instantiation in a truncated hydrogen
basis, and sketch a Lean~4 formalisation strategy treating $\Xi(t)$ as
a sheaf morphism over the combined number/physics topos.
\end{abstract}
\newpage
\tableofcontents
\newpage
\section{Executive Summary}

This report consolidates three strands of work:
\begin{enumerate}[label=\textbf{(\alph*)}]
  \item The Hydrogen Universe Standard (H-US), which introduces
  hydrogen-based units (Hydrogen Radius $H_R$, Hydrogen Tick $H_T$,
  Hydrogen Quantum $H_Q$) and a prime-indexed recursive tensor
  framework for physical quantities.

  \item The Alpha-Hydrogen programme, which deforms standard hydrogen
  wavefunctions by an ``alpha'' factor $\alpha_n(\rho)$ to encode
  multiplicity effects, especially in the radial sector of the
  Schr\"odinger solution.

  \item The $\Phi$-Quantum Natural Number Singularity ($\Phi$-QNNS),
  where $\mathbb{N}$, $\mathbb{Q}$, $\mathbb{R}$, and $\mathbb{C}$
  emerge as gauge-invariant or attractor states of a prime-indexed,
  zeta-modulated tensor evolution, realised as a stratified sheaf in a
  $\Phi$-topos.
\end{enumerate}

The central advancement is a \emph{unified prime-weighted evolution
operator} $\Xi(t)$ with the following features:
\begin{itemize}
  \item A shared prime frequency $\omega_p = \log p$ for each prime
  $p$, appearing both in the hydrogenic and number-field dynamics.

  \item A zeta-modulated multiplicity field
  \[
  \Lambda_m(p,t)
    = \Lambda_m^{(0)}(p)\,\Phi\,p^{-\alpha}
      \,\Re\bigl[\zeta\bigl(\tfrac12 + i\,\omega_p t\bigr)\bigr],
  \]
  which couples primes, time $t$, and the golden ratio $\Phi$.

  \item A non-commutative ordered product
  \[
  U_{\Xi}(t)
    = \prod_{p\le p_{\max}}
        \exp\bigl(-i\,w_p(t)\,\Lambda_m(p,t)\,\mathcal{D}_p\bigr),
  \]
  where $\mathcal{D}_p$ is a prime-indexed operator acting on the
  relevant fiber (hydrogenic Hilbert space or number sheaf), and where
  $w_p(t)$ contains a log-periodic drive.
\end{itemize}

On the hydrogen side, $U_{\Xi}(t)$ acts in a truncated $s$-state
basis to produce an evolved state $\psi(t)$ whose probability density
in $\log\rho$ exhibits enhanced power in frequency bands near
$\omega_p = \log p$. On the $\Phi$-QNNS side, the same prime weights
enter a covariant derivative whose vanishing picks out natural numbers
and drives the extension to $\mathbb{Q}$, $\mathbb{R}$, and
$\mathbb{C}$.

Finally, we propose a Lean~4 formalisation route: define a stratified
sheaf of fibers (number objects and H-US Hilbert sector), attach a
prime-indexed family of operators, and encode $U_{\Xi}(t)$ as an
ordered fold/product. This provides a shared logical substrate for
both the physical and number-theoretic realizations of $\Xi(t)$.

\section{Conceptual Overview}

\subsection{Hydrogen Universe Standard and Alpha-Hydrogen}

The Hydrogen Universe Standard (H-US) takes hydrogen as a metrological
foundation, defining:
\begin{align*}
  H_R &:= \text{Bohr radius} \approx 0.529\times10^{-10}\,\mathrm{m},\\
  H_T &:= \text{Rydberg period} \approx 4.14\times10^{-16}\,\mathrm{s},\\
  H_Q &:= \text{ground-state energy} \approx 13.6\,\mathrm{eV}.
\end{align*}
These units induce dimensionless hydrogenic coordinates, such as
\(\rho = r / H_R\), and support Lorentz-covariant transformations
$T_{H\to\mathrm{SI}}$ via scalar factors $\kappa_X$.

The Alpha-Hydrogen programme introduces an extra factor
$\alpha_n(\rho)$ into the standard hydrogen radial solutions. For
$l=0$ states, one may write
\[
  R_{n0}(\rho)
    \propto e^{-\rho/n}
    L^{1}_{n-1}\!\left(\frac{2\rho}{n}\right),
\]
and consider deformations
\[
  \tilde R_{n0}(\rho,t)
    = \alpha_n(\rho;t)\,R_{n0}(\rho),
\]
where $\alpha_n$ is intended to encode prime-indexed multiplicity and
recursive structure.

\subsection{The $\Phi$-Quantum Natural Number Singularity}

The $\Phi$-QNNS framework treats natural numbers as ground states of a
prime-indexed $U(1)$ gauge field. For each prime $p$, one defines a
gauge mode with frequency $\omega_p = \log p$ and a zeta-modulated
amplitude. A typical gauge field term is
\[
  A_p(t)
    = \Lambda_m\,p^{-\alpha}\,e^{i\omega_p t},
\]
with $\alpha > 1$, and the tensor $T$ obeys a covariant derivative
equation schematically of the form
\[
  D_t T = \frac{dT}{dt} - i \sum_p A_p(t)\,(\cdots)\,T.
\]
Zeros of $D_t T$ occur at integer $t \in \mathbb{N}$ for simple
choices of $T(t)$, and the framework extends to rationals,
reals, and complex numbers via $p$-adic completions, adelic flows, and
modular/monodromy structures.

The number hierarchy
\[
  \mathbb{N} \subset \mathbb{Q} \subset \mathbb{R} \subset \mathbb{C}
\]
is then viewed as a stratified sheaf in a $\Phi$-topos, with primes
serving as ``neurons'' in a cosmic neural network, and $\zeta$-related
weights setting coupling strengths.

\subsection{Unified Vision}

The key observation is that both H-US/Alpha-Hydrogen and $\Phi$-QNNS
use:
\begin{itemize}
  \item Prime frequencies $\omega_p = \log p$,
  \item A multiplicity constant or field $\Lambda_m$,
  \item A recursive, time-dependent operator $\Xi(t)$ acting on
  tensor-like objects,
  \item Sheaf/topos language for gluing local structures into a global
  object.
\end{itemize}
We therefore introduce a single evolution operator $\Xi(t)$ whose
action has two principal specializations:
\begin{enumerate}[label=\textbf{(\roman*)}]
  \item On the H-US hydrogen Hilbert space, $\Xi(t)$ deforms the
  hydrogen $s$-state basis via a non-commutative product of prime
  dilation operators, giving log-radius spectra with prime-enhanced
  band power.

  \item On the number-sheaf fibers, $\Xi(t)$ induces a gauge-like
  recursion whose fixed points and attractors realise the number
  hierarchy and $\Phi$-QNNS dynamics.
\end{enumerate}

\section{Unified Prime-Weighted Evolution Operator \texorpdfstring{$\Xi(t)$}{}}

\subsection{Shared Prime Field and Weights}

Let $P$ be the set of primes. For each $p \in P$, define the prime
frequency
\[
  \omega_p := \log p.
\]
We introduce a \emph{zeta-modulated multiplicity field}
\[
  \Lambda_m(p,t)
    := \Lambda_m^{(0)}(p)\,\Phi\,p^{-\alpha}
      \,\Re\bigl[\zeta\bigl(\tfrac12 + i\,\omega_p t\bigr)\bigr],
\]
where:
\begin{itemize}
  \item $\Phi = \frac{1+\sqrt{5}}{2}$ is the golden ratio;
  \item $\alpha > 1$ is a decay exponent ensuring convergence of prime
  sums;
  \item $\Lambda_m^{(0)}(p)$ is a static prime weight, such as
  $\log p / \log p_{\max}$ in H-US/PIRTM.
\end{itemize}
We also define a log-periodic drive amplitude
\[
  w_p(t) := \frac{\log p}{p}\,\sin(\omega_p t).
\]

\subsection{Abstract Operator Form}

Let $\mathcal{S}$ be a sheaf (or sheaf-like structure) whose stalks
are:
\begin{itemize}
  \item Number objects: $\mathbb{N}$, $\mathbb{Q}$, $\mathbb{R}$,
  $\mathbb{C}$,
  \item The H-US hydrogenic Hilbert spaces (e.g.\ radial $s$-state
  Hilbert spaces in H-US units).
\end{itemize}
For each prime $p$, we posit a self-adjoint operator $\mathcal{D}_p$
acting on the stalk at hand:
\begin{itemize}
  \item On number stalks: $\mathcal{D}_p^{\text{num}}$ is the prime
  gauge/tensor generator entering the $\Phi$-QNNS covariant derivative.

  \item On hydrogenic H-US stalks: $\mathcal{D}_p^{\text{phys}}$ is
  the prime-specific dilation/log-periodic operator (matrix $D_p$ in a
  truncated hydrogen basis).
\end{itemize}
We write $\mathcal{D}_p$ generically, with the understanding that it
restricts appropriately to each fiber.

The unified evolution operator is then defined for a finite prime
cutoff $p_{\max}$ by the ordered product
\begin{equation}
  U_{\Xi}(t)
    := \prod_{p\in P_{\le p_{\max}}}
        \exp\bigl(-i\,w_p(t)\,\Lambda_m(p,t)\,\mathcal{D}_p\bigr).
  \label{eq:U_Xi_def}
\end{equation}
Here:
\begin{itemize}
  \item The product order over primes is fixed (e.g.\ ascending $p$),
  making $U_\Xi(t)$ a non-commutative product.
  \item On a given stalk/fiber $F$, $U_\Xi(t)$ is a (formal) unitary
  when each $\mathcal{D}_p$ is self-adjoint.
\end{itemize}

\subsection{Specializations}

\subsubsection{Number-Field / $\Phi$-QNNS Specialization}

On the number-sheaf stalks, $\Xi(t)$ induces a covariant derivative
\[
  D_t T
    = \frac{dT}{dt} - i\sum_{p} A_p(t)\,T,
\]
where 
\[
  A_p(t)
    \propto \Lambda_m(p,t)\,e^{i\omega_p t}.
\]
In the $\Phi$-QNNS setting, one finds that $D_t T = 0$ picks out
integer times $t \in \mathbb{N}$ for certain simple test tensors (e.g.
$T(t) = t^2$), and further constructions lift this to $\mathbb{Q}$,
$\mathbb{R}$, and $\mathbb{C}$ via $p$-adic, adelic, and modular
structures.

\subsubsection{Hydrogenic / H-US Specialization}

On the H-US hydrogenic Hilbert space (e.g.\ truncated $s$-state basis
$n=1,\dots,N$), $\mathcal{D}_p$ is realised numerically as a matrix
combining a dilation operator with a $\cos(\log_p \rho)$ potential:
\[
  \mathcal{D}_p^{\text{phys}}
    = \hat D_{\text{dilate}} + \epsilon \cos(\log_p \rho),
\]
where $\hat D_{\text{dilate}}$ corresponds (up to factors) to
$\rho\,\partial_\rho + \tfrac12$, and $\epsilon$ controls the strength
of the log-periodic modulation. Then $U_\Xi(t)$ acts on an initial
state $\psi(0)$ to yield $\psi(t) = U_\Xi(t)\psi(0)$. In particular,
for $\psi(0)$ taken as the ground $1s$ state, one can study the
probability density $|\psi(\rho,t)|^2$ in $\log\rho$ and its Fourier
spectrum.

Numerical experiments (not reported here in detail) indicate:
\begin{itemize}
  \item Overlaps $|\langle 1s | \psi(t)\rangle|^2$ in a high but
  nontrivial range (e.g.\ $\sim 0.85$--$0.95$), confirming that
  $U_\Xi(t)$ is a structured deformation rather than a randomizer.

  \item Enhanced power in frequency bands around $\omega = \log p$
  compared to nearby non-prime control frequencies, as measured by
  band power ratios and permutation tests.
\end{itemize}
These hydrogenic spectral features are then interpreted as a physical
manifestation of the same prime-harmonic structure governing the
number-field recursion.

\subsection{Sheaf-Morphism Perspective}

Let $\mathcal{S}$ denote the stratified sheaf whose fibers are the
number objects and the H-US hydrogenic Hilbert spaces. Then:
\[
  \Xi(t) : \mathcal{S} \longrightarrow \mathcal{S}
\]
is a sheaf morphism whose local action on each stalk is given by the
restriction of $U_\Xi(t)$ in Eq.~\eqref{eq:U_Xi_def} to that fiber.
The unification is thus conceptual as well as operational: there is a
single prime-weighted engine $\Xi(t)$ which, when viewed on different
fibers, realises both number and spectral dynamics.

\section{Numerical Implementation for Hydrogenic States}

\subsection{Truncated Hydrogen Basis}

We consider hydrogen $s$-states ($l=0$) with principal quantum number
$n=1,\dots,N$ and radial coordinate $r$ (with $H_R$ set to $1$ for
numerical convenience). The standard normalized radial wavefunctions
$R_{n0}(r)$ are given by
\[
  R_{n0}(r)
    = \mathcal{N}_n\,e^{-\rho/2} L_{n-1}^{1}(\rho),
  \quad \rho = \frac{2r}{n},
\]
with $L_{n-1}^{1}$ an associated Laguerre polynomial and
$\mathcal{N}_n$ a normalization constant. For numerical quadrature,
one works with
\[
  \Phi_n(r) := R_{n0}(r)\,r,
\]
so that the radial measure becomes $\mathrm{d}r$ (the $r^2$ factor is
absorbed).

A suitable radial grid $r_i$ (e.g.\ uniform from $10^{-6}$ to $r_{\max}$)
and weights $w_i$ are chosen to approximate integrals via
\[
  \int_0^{r_{\max}} f(r)\,r^2\,\mathrm{d}r
    \approx \sum_i w_i\,f(r_i).
\]

\subsection{Construction of \texorpdfstring{$\mathcal{D}^{\text{phys}}_p$}{}}

The dilation operator is approximated numerically by:
\begin{align*}
  \left(\hat D_{\text{dilate}} \Phi_j\right)(r)
    &:= -i\left(r\,\partial_r + \tfrac12\right)\Phi_j(r),\\
  \left(\mathcal{D}^{\text{phys}}_p \Phi_j\right)(r)
    &:= \hat D_{\text{dilate}}\Phi_j(r) + \epsilon \cos(\log_p r)\,\Phi_j(r),
\end{align*}
with $\epsilon$ a small parameter controlling log-periodic driving.
Matrix elements in the basis $\{\Phi_n\}_{n=1}^N$ are computed via
\[
  (D_p)_{ij}
    = \langle i | \mathcal{D}^{\text{phys}}_p | j \rangle
    \approx \sum_k w_k\,\Phi_i(r_k)\,
        \left(\mathcal{D}^{\text{phys}}_p\Phi_j\right)(r_k),
\]
yielding an $N\times N$ Hermitian matrix $D_p$.

\subsection{Evolution and Diagnostics}

Given a list of primes $p \in \{2,3,5,\dots,p_{\max}\}$ and the
weights $w_p(t)\Lambda_m(p,t)$, the evolution operator restricted to
hydrogenic $s$-states is
\[
  U_{\Xi}^{\text{phys}}(t)
    = \prod_{p\le p_{\max}} \exp\bigl(-i\,w_p(t)\,\Lambda_m(p,t)\, D_p\bigr),
\]
with the product taken in ascending $p$. For the initial vector
$\psi(0) = (1,0,\dots,0)$ corresponding to the $1s$ state, we define
\[
  \psi(t) := U_{\Xi}^{\text{phys}}(t)\,\psi(0).
\]

Reconstructing the radial wavefunction on the grid via
\[
  \psi(r_k,t) \approx \sum_{n=1}^{N} \psi_n(t)\,\Phi_n(r_k)/r_k,
\]
one can study:
\begin{enumerate}[label=\textbf{(\alph*)}]
  \item The overlap with the initial state,
  \[
    \text{overlap}(t)
      := |\langle 1s | \psi(t)\rangle|^2
      = |\psi_1(t)|^2.
  \]
  \item The probability density in $\log r$, obtained by interpolating
  $|\psi(r,t)|^2$ onto a uniform $\log r$ grid and computing the FFT
  spectrum.

  \item Band powers around target frequencies $\omega_p = \log p$ and
  near-by control frequencies, to test for enhanced prime-band power.
\end{enumerate}

Permutation tests can then be used to assess the statistical
significance of observed prime vs.\ control band power differences.

\section{Illustrative Python Snippet}

The following is a schematic Python snippet (omitting some details)
that captures the hydrogenic specialization of $U_\Xi(t)$ with
zeta-modulated weights:
\begin{verbatim}
import numpy as np
from scipy.linalg import expm
from mpmath import zeta, re

def Lambda_m(p, t, p_max, alpha=1.6):
    phi = (1 + np.sqrt(5.0)) / 2.0
    Lambda0 = np.log(p) / np.log(p_max)
    zeta_val = float(re(zeta(0.5 + 1j * np.log(p) * t)))
    return Lambda0 * phi * (p ** (-alpha)) * zeta_val

def w_p_times_Lambda(p, t, p_max):
    omega = np.log(p)
    w = (np.log(p) / p) * np.sin(omega * t)
    return w * Lambda_m(p, t, p_max)

def evolve_state(psi0, primes, t, Dp_matrices, p_max):
    psi = psi0.copy()
    for p, Dp in zip(primes, Dp_matrices):
        coeff = w_p_times_Lambda(p, t, p_max)
        U = expm(-1j * coeff * Dp)
        psi = U @ psi
    return psi
\end{verbatim}

Here, \verb|Dp_matrices| is a list of $D_p$ matrices constructed as
described earlier, \verb|psi0| is the initial state, and
\verb|primes| is the finite prime set. The logarithmic frequency
analysis and permutation tests are implemented on top of the
reconstructed $\psi(r,t)$ data.

\section{Lean 4 Formalisation Sketch}

\subsection{Fibers and Sheaf Structure}

We introduce an inductive type representing the different fibers:
\[
  \texttt{XiFiber} :=
    \{\texttt{natF}, \texttt{ratF}, \texttt{realF}, \texttt{complexF},
    \texttt{husHilbert}\}.
\]

In Lean-like pseudocode:
\begin{verbatim}
inductive XiFiber
| natF | ratF | realF | complexF | husHilbert
deriving DecidableEq, Repr

structure XiSheaf where
  Section : XiFiber → Type
  restrict :
    {a b : XiFiber} → a = b → Section a → Section b
\end{verbatim}

\subsection{Prime Operators and Weights}

We encode prime-indexed operators and multiplicity weights:

\begin{verbatim}
structure PrimeOp (S : XiSheaf) where
  op : Nat → {f : XiFiber} → S.Section f → S.Section f
  prime_only : ∀ n, Nat.Prime n → True

structure XiWeights where
  Lambda : Nat → Real → Real
\end{verbatim}

The unified $\Xi$-data:
\begin{verbatim}
structure UnifiedXi (S : XiSheaf) where
  D : PrimeOp S
  W : XiWeights
\end{verbatim}

\subsection{Ordered Evolution}

A finite prime list encodes the cutoff $p_{\max}$. We define an
ordered action on any given fiber:

\begin{verbatim}
def orderedXi
  (S : XiSheaf)
  (U : UnifiedXi S)
  (primes : List Nat)
  (t : Real) :
  {f : XiFiber} → S.Section f → S.Section f :=
by
  intro f x
  exact
    primes.foldl
      (fun acc p => U.D.op p acc)
      x
\end{verbatim}

This definition abstracts away analytic details (such as the
exponentials and unitarity) and captures the key structural fact: a
prime-ordered evolution built from prime-indexed operators. The
hydrogenic and number-field specializations correspond to choosing
concrete types for \verb|S.Section husHilbert| and \verb|S.Section natF|
etc., together with interpretations of \verb|U.D.op p|.

\subsection{Initial Theorem Targets}

Natural early lemmas and theorems include:
\begin{itemize}
  \item \textbf{Identity at empty cutoff:} if \verb|primes = []|,
  then \verb|orderedXi S U [] t| is the identity on each fiber.

  \item \textbf{Fiber preservation:} \verb|orderedXi| maps each fiber
  \verb|f| to itself (no cross-fiber mixing without additional
  structure).

  \item \textbf{Functoriality in prime lists:} extending the prime list
  corresponds to composing the corresponding operators, preserving the
  order.

  \item \textbf{Existence of specializations:} there exist instances of
  \verb|XiSheaf| and \verb|UnifiedXi| corresponding to (a) the H-US
  hydrogenic Hilbert sector and (b) the $\Phi$-QNNS number sheaf.
\end{itemize}
These Lean-level statements give a formally checked backbone for the
informal mathematics presented earlier.

\section{Conclusion and Outlook}

We have:
\begin{itemize}
  \item Defined a unified prime-weighted evolution operator $\Xi(t)$
  that acts as a sheaf morphism across both hydrogenic and
  number-theoretic fibers.

  \item Shown how this operator specialises to a matrix evolution on a
  truncated hydrogen basis, delivering observable log-radius spectral
  features at prime frequencies.

  \item Sketched how the same prime weights and frequencies govern the
  $\Phi$-QNNS recursion on number-sheaf layers, yielding natural numbers
  and higher number systems as fixed points and attractors.

  \item Outlined a Lean~4 encoding of fibers, prime operators, weights,
  and ordered prime evolution, paving the way for machine-checked
  reasoning about $\Xi(t)$.
\end{itemize}

Next steps include:
\begin{enumerate}
  \item A full parameter scan and statistical analysis of the hydrogenic
  spectral signatures under $\Xi(t)$, especially with zeta-modulated
  weights.

  \item A more detailed categorical formulation of the sheaf $\mathcal{S}$
  and the precise conditions under which $\Xi(t)$ is a natural
  transformation.

  \item Comprehensive Lean~4 development, tying together H-US primed
  tensors, $\Phi$-QNNS number sheaves, and the unified $\Xi(t)$
  framework.
\end{enumerate}

Taken together, these components support the view that multiplicity,
primes, and log-periodicity provide a common engine for both physical
spectra and the emergence of number systems.

\nocite{*}
\bibliographystyle{plainnat}
\bibliography{references} 
\end{document}